\chapter{Términos de Referencia (TdR)}

\section{Objeto de la Contratación}
Contratar la ejecución de la obra para el montaje, instalación, pruebas y puesta en servicio de las **Instalaciones Eléctricas Industriales de la Nave Industrial de 20x40m**, garantizando el cumplimiento del Código Nacional de Electricidad - Utilización, las Normas RNE EM.010 y las especificaciones técnicas del presente expediente.

\section{Perfil y Experiencia del Contratista}
\begin{itemize}
    \item El contratista o ejecutor deberá contar con registro vigente de proveedor de obras eléctricas.
    \item El Ingeniero Residente de Obra deberá ser un **Ingeniero Mecánico Eléctrico o Eléctrico Colegiado y Habilitado (CIP)**, con no menos de 3 años de experiencia en montaje de plantas e instalaciones industriales.
\end{itemize}

\section{Obligaciones y Normas de Seguridad}
1. Cumplimiento obligatorio de la **Norma G.050 Seguridad durante la Construcción** y el uso permanente de Equipos de Protección Personal (EPP) de tipo dieléctrico (casco, botas $10\text{ kV}$, guantes dieléctricos y arnés de seguridad para trabajos en altura sobre $1.80\text{ m}$).
2. Tramitación de permisos y coordinación de la interrupción del suministro con la Empresa Concesionaria de Distribución Eléctrica local (Electro Puno S.A.A.).

\section{Pruebas, Recepción y Garantía de Obra}
1. El contratista entregará los protocolos de prueba firmados por el Ingeniero Residente:
   - Protocolo de resistencia de aislamiento de todos los conductores N2XH ($\ge 50 \text{ M}\Omega$).
   - Protocolo de medición de resistencia de la malla de puesta a tierra ($\le 5.0 \:\Omega$).
   - Protocolo de tiempos de disparo de protecciones diferenciales SI ($\le 30 \text{ ms}$).
   - Protocolo de secuencia de fases y equilibrio de cargas en TGD/TF1.
2. El plazo de garantía de los equipos y la instalación será de no menos de **12 meses** contados a partir del acta de Recepción de Obra.
